% Bagean: metodhe
% Bab: bukti
% Pérangan: maca-bukti

\documentclass[../../../include/open-logic-section]{subfiles}

\begin{document}

\olfileid{mth}{prf}{rea}
\olsection{Maca Bukti}

Bukti ing buku ajar lan artikel arang banget ngemot kabèh
rincian kaya ing conto-conto sadurungé. Pangarang kerep ora
nyebut nalika mbédakaké kasus, nalika nulis bukti ora langsung,
utawa nalika migunakaké dhéfinisi. Dadi nalika maca bukti ing
buku ajar, panjenengan asring kudu nglengkapi rincian mau
dhéwé supaya bisa ngerti buktiné. Iki uga latihan kang apik
kanggo nguwasani manéka langkah ing bukti. Ayo deleng conto.

\begin{prop}[Absorpsi]
Kanggo kabèh himpunan $A$, $B$,
\[
A \cap (A \cup B) = A
\]
\end{prop}

\begin{proof}
Yèn $z \in A \cap (A \cup B)$, mula $z \in A$, saéngga
$A \cap (A \cup B) \subseteq A$. Saiki anggepen $z \in A$.
Mula uga $z \in A \cup B$, lan mulané
$z \in A \cap (A \cup B)$.
\end{proof}

Bukti hukum absorpsi mau cekak banget. Ora ana dhéfinisi
kang disebut, ora ana ukara ``kang kudu dibuktèkaké yaiku''
sadurungé mbuktekaké, lan sapanunggalané. Ayo dijlentrehaké.
Proposisi mau pratelan umum kanggo sembarang himpunan $A$ lan
$B$, lan nalika bukti nyebut $A$ utawa $B$, iki variabel
kanggo himpunan sembarang. Pratelan umum kang dibuktèkaké
yaiku kang dibutuhaké kanggo pepadhan himpunan: saben
!!{element} ing sisih kiwa pepadhan uga !!a{element}
ing sisih tengen, lan kosok baliné.

\begin{quote}
``Yèn $z \in A \cap (A \cup B)$, mula $z \in A$, saéngga
  $A \cap (A \cup B) \subseteq A$.''
\end{quote}

Iki separo kapisan saka bukti pepadhan: nuduhaké yèn
sebarang~$z$ kang !!a{element} sisih kiwa uga
!!a{element} sisih tengen, yaiku
$A \cap (A \cup B) \subseteq A$.
Anggepen $z \in A \cap (A \cup B)$. Amarga $z$
!!{element} irisan rong himpunan yèn lan mung yèn
iku !!{element} loro-loroné, kita nyimpulaké
$z \in A$ lan uga $z \in A \cup B$.
Mliginé, $z \in A$, kang pancèn arep dibuktèkaké.
Amarga separo kapisan mung mbutuhaké iki, sisane
bukti mesthi kanggo separo kapindho, yaiku
$A \subseteq A \cap
(A \cup B)$.

\begin{quote}
``Saiki anggepen $z \in A$. Mula uga $z \in A \cup B$, lan
mulané $z \in A \cap (A \cup B)$.''
\end{quote}

Kita miwiti kanthi nganggep $z \in A$, amarga arep
nuduhaké yèn kanggo sembarang~$z$, yèn $z \in A$
mula $z \in A \cap (A \cup B)$. Kanggo mbuktekaké
$z \in A \cap (A \cup B)$, miturut dhéfinisi
``$\cap$'', kita kudu nuduhaké (i) $z \in A$
lan (ii) $z \in A \cup B$. Pérangan (i) wis
dadi anggepan, mula ora perlu dibuktèkaké maneh
lan ora disebut maneh ing bukti. Kanggo (ii), élinga
yèn $z$ iku !!a{element} gabungan rong himpunan
yèn lan mung yèn iku !!{element} paling ora salah
sijiné. Amarga $z \in A$, lan $A \cup B$ iku
gabungan $A$ lan $B$, syarat iki kacukupan.
Dadi $z \in A \cup B$. Kita wis nuduhaké
(i) $z \in A$ lan (ii) $z \in A \cup B$;
mula miturut dhéfinisi ``$\cap$'',
$z \in A \cap (A \cup B)$. Bukti ringkes
ora nyebut dhéfinisi mau amarga dianggep pamaca
wis mangertèni. Yèn durung, panjenengan
kudu ngélingi maneh dhéfinisiné. Banjur
panjenengan uga kudu ngerti ngapa saka $z
\in A$ katut $z \in A \cup B$, lan saka
$z \in A$ lan $z \in A \cup B$ katut
$z \in A \cap (A \cup B)$.

Iki versi liya saka bukti mau kang njlentrehaké
kabèh langkah:
\begin{proof}{}
[Miturut dhéfinisi $=$ kanggo himpunan, kanggo
  mbuktekaké $A \cap (A \cup B) = A$
  kita kudu nuduhaké (a)
  $A \cap (A \cup B) \subseteq A$ lan (b)
  $A \subseteq A \cap (A \cup B)$.
  (a): Miturut dhéfinisi $\subseteq$, kita kudu
  nuduhaké yèn $z \in A \cap (A \cup B)$,
  mula $z \in A$.] Yèn $z \in A
\cap (A \cup B)$, mula $z \in A$
[amarga miturut dhéfinisi $\cap$,
$z \in A \cap (A \cup B)$ yèn lan mung
yèn $z \in A$ lan $z \in A \cup B$],
mula $A \cap (A \cup B) \subseteq A$.
[(b): Miturut dhéfinisi $\subseteq$,
kita kudu nuduhaké yèn $z \in A$,
mula $z \in A \cap (A \cup B)$.]
Saiki anggepen [(1)] $z \in A$.
Mula uga [(2)] $z \in A \cup B$
[amarga miturut (1) $z \in A$ utawa
$z \in B$, kang miturut dhéfinisi $\cup$
ateges $z \in A \cup B$], lan mulané
$z \in A \cap (A \cup B)$
[amarga dhéfinisi $\cap$ mbutuhaké
$z \in A$, yaiku (1), lan
$z \in A \cup B$, yaiku (2)].
\end{proof}

\begin{prob}
Jlentrehna bukti ing ngisor iki kanggo
$A \cup (A \cap B) = A$, kanthi nyebut kabèh
pola inferensi kang digunakaké, sebab saben
langkah katut saka anggepan utawa pratelan kang
wis kabukten, lan dhéfinisi endi kang ditrapaké.
\begin{proof}
  Yèn $z \in A \cup (A \cap B)$ mula $z \in A$
  utawa $z \in A \cap B$. Yèn
  $z \in A \cap B$, mula $z \in A$.
  Saben $z \in A$ uga $z \in A \cup (A
  \cap B)$.
\end{proof}
\end{prob}

\end{document}
